% Options for packages loaded elsewhere
\PassOptionsToPackage{unicode}{hyperref}
\PassOptionsToPackage{hyphens}{url}
%
\documentclass[
]{article}
\usepackage{amsmath,amssymb}
\usepackage{iftex}
\ifPDFTeX
  \usepackage[T1]{fontenc}
  \usepackage[utf8]{inputenc}
  \usepackage{textcomp} % provide euro and other symbols
\else % if luatex or xetex
  \usepackage{unicode-math} % this also loads fontspec
  \defaultfontfeatures{Scale=MatchLowercase}
  \defaultfontfeatures[\rmfamily]{Ligatures=TeX,Scale=1}
\fi
\usepackage{lmodern}
\ifPDFTeX\else
  % xetex/luatex font selection
\fi
% Use upquote if available, for straight quotes in verbatim environments
\IfFileExists{upquote.sty}{\usepackage{upquote}}{}
\IfFileExists{microtype.sty}{% use microtype if available
  \usepackage[]{microtype}
  \UseMicrotypeSet[protrusion]{basicmath} % disable protrusion for tt fonts
}{}
\makeatletter
\@ifundefined{KOMAClassName}{% if non-KOMA class
  \IfFileExists{parskip.sty}{%
    \usepackage{parskip}
  }{% else
    \setlength{\parindent}{0pt}
    \setlength{\parskip}{6pt plus 2pt minus 1pt}}
}{% if KOMA class
  \KOMAoptions{parskip=half}}
\makeatother
\usepackage{xcolor}
\usepackage{longtable,booktabs,array}
\usepackage{calc} % for calculating minipage widths
% Correct order of tables after \paragraph or \subparagraph
\usepackage{etoolbox}
\makeatletter
\patchcmd\longtable{\par}{\if@noskipsec\mbox{}\fi\par}{}{}
\makeatother
% Allow footnotes in longtable head/foot
\IfFileExists{footnotehyper.sty}{\usepackage{footnotehyper}}{\usepackage{footnote}}
\makesavenoteenv{longtable}
\setlength{\emergencystretch}{3em} % prevent overfull lines
\providecommand{\tightlist}{%
  \setlength{\itemsep}{0pt}\setlength{\parskip}{0pt}}
\setcounter{secnumdepth}{-\maxdimen} % remove section numbering
% definitions for citeproc citations
\NewDocumentCommand\citeproctext{}{}
\NewDocumentCommand\citeproc{mm}{%
  \begingroup\def\citeproctext{#2}\cite{#1}\endgroup}
\makeatletter
 % allow citations to break across lines
 \let\@cite@ofmt\@firstofone
 % avoid brackets around text for \cite:
 \def\@biblabel#1{}
 \def\@cite#1#2{{#1\if@tempswa , #2\fi}}
\makeatother
\newlength{\cslhangindent}
\setlength{\cslhangindent}{1.5em}
\newlength{\csllabelwidth}
\setlength{\csllabelwidth}{3em}
\newenvironment{CSLReferences}[2] % #1 hanging-indent, #2 entry-spacing
 {\begin{list}{}{%
  \setlength{\itemindent}{0pt}
  \setlength{\leftmargin}{0pt}
  \setlength{\parsep}{0pt}
  % turn on hanging indent if param 1 is 1
  \ifodd #1
   \setlength{\leftmargin}{\cslhangindent}
   \setlength{\itemindent}{-1\cslhangindent}
  \fi
  % set entry spacing
  \setlength{\itemsep}{#2\baselineskip}}}
 {\end{list}}
\usepackage{calc}
\newcommand{\CSLBlock}[1]{\hfill\break\parbox[t]{\linewidth}{\strut\ignorespaces#1\strut}}
\newcommand{\CSLLeftMargin}[1]{\parbox[t]{\csllabelwidth}{\strut#1\strut}}
\newcommand{\CSLRightInline}[1]{\parbox[t]{\linewidth - \csllabelwidth}{\strut#1\strut}}
\newcommand{\CSLIndent}[1]{\hspace{\cslhangindent}#1}
\ifLuaTeX
  \usepackage{selnolig}  % disable illegal ligatures
\fi
\usepackage{bookmark}
\IfFileExists{xurl.sty}{\usepackage{xurl}}{} % add URL line breaks if available
\urlstyle{same}
\hypersetup{
  hidelinks,
  pdfcreator={LaTeX via pandoc}}

\author{}
\date{}

\begin{document}

\section{Temporal Evaluation, Calibrated Risk, and Conformal
Delay-Severity Prediction in Pharmaceutical
Logistics}\label{temporal-evaluation-calibrated-risk-and-conformal-delay-severity-prediction-in-pharmaceutical-logistics}

\textbf{Anonymized manuscript for double-anonymized review}

\subsection{Abstract}\label{abstract}

Delivery delays in global health supply chains disrupt treatment
continuity, deplete buffer stocks, and complicate procurement and
transportation planning. Machine-learning models for shipment-delay
prediction are often evaluated with random cross-validation, which can
obscure deterioration when the data-generating process changes over
time. We evaluate an integrated logistics decision-support framework
that combines prediction-time leakage control, expanding-origin temporal
evaluation with 90-day post-delivery embargoes, calibrated delay-risk
estimation, conditional delay-severity modeling, conformalized quantile
uncertainty, and capacity-constrained shipment prioritization. Using the
USAID Supply Chain Management System delivery-history data (N = 10,324
shipments across 42 recipient countries), random cross-validation
overestimated PR-AUC relative to temporally ordered evaluation across
all five tested classifiers, with relative inflation from 26.2\% for
logistic regression to 99.7\% for LightGBM. Platt scaling reduced Brier
score and expected calibration error while preserving continuous
rankings. Conditional-median severity estimates achieved lower mean
absolute error than LightGBM quantile point predictions, while the
quantile models enabled Split Conformalized Quantile Regression. Under
exchangeability, split CQR has finite-sample marginal coverage
guarantees; here, temporally later cohorts are treated as empirical
out-of-time coverage audits rather than formal guarantees under shift.
On the secondary locked registry benchmark (N = 1,013; 61 delays), 90\%
CQR achieved 91.80\% empirical coverage (56/61; exact 95\% CI
81.90\%-97.28\%) with 46.27-day mean width. Simulated Top-K triage
showed different trade-offs among risk, severity, and uncertainty. The
results support temporally disciplined and uncertainty-aware logistics
evaluation rather than a single universally dominant model or ranking
policy.

\textbf{Keywords:} Pharmaceutical logistics; shipment delay; temporal
validation; probability calibration; conformal prediction; decision
support

\subsection{1. Introduction}\label{introduction}

Global health supply chains operate under long international lead times,
heterogeneous transport modes, procurement constraints, and uneven
infrastructure. Delayed health-product deliveries can contribute to
stockouts and emergency responses, making delivery reliability a
practical logistics concern (Yadav 2015; Vledder et al. 2019).
Predictive analytics can help identify shipments that warrant attention,
but a deployable logistics model must answer more than whether an order
is likely to be late. It must be evaluated at a realistic prediction
time, provide probabilities that are meaningful for decision thresholds,
characterize how severe a delay may become, communicate uncertainty, and
support prioritization when review capacity is limited.

A central methodological difficulty is temporal dependence. Randomly
mixing historical and future shipments can allow a model-selection
process to exploit distributional information unavailable at deployment
time. Leakage and poorly matched validation can therefore produce
optimistic scientific claims (Kapoor and Narayanan 2023; Roberts et al.
2017; Bergmeir and Benítez 2012). In shipment data, the problem is
amplified by label maturity: an order can be known before its eventual
delivery outcome is resolved. We therefore treat prediction-time
integrity and temporal separation as part of the scientific problem
rather than as implementation details.

A second challenge is that ranking and probability reliability are
different objectives. A classifier may separate late and on-time
shipments reasonably well while assigning probability values that are
unsuitable for thresholding or resource allocation. Post-hoc calibration
methods such as Platt scaling and isotonic regression provide
established tools for this distinction (Platt 1999; Niculescu-Mizil and
Caruana 2005; Guo et al. 2017). Third, binary risk and delay magnitude
are operationally distinct. A short delay and a prolonged disruption can
have different consequences even when both share the same binary label.
Quantile regression offers a way to describe asymmetric conditional
delay distributions (Koenker and Bassett 1978), while conformalized
quantile regression (CQR) can conformalize such bounds under its
validity assumptions (Romano, Patterson, and Candès 2019; Angelopoulos
and Bates 2023).

Finally, logistics teams face finite intervention capacity. Prediction
becomes operationally useful only when it informs which shipments to
inspect or escalate. This motivates a decision layer in which calibrated
risk, conditional severity, and uncertainty may produce different
rankings rather than a claim that one score is universally superior
(Bertsimas and Kallus 2020).

This study evaluates these components within historical USAID Supply
Chain Management System (SCMS) shipment records. Its contributions are:
(1) a leakage-aware expanding temporal evaluation with a 90-day
label-maturity embargo; (2) a five-model comparison quantifying
random-split optimism; (3) separate evaluation of probability
calibration and conditional severity; (4) empirical auditing of CQR
coverage and sharpness in temporally later cohorts; and (5) a
capacity-constrained Top-K prioritization analysis explicitly labeled as
a simulated operational scenario.

\subsection{2. Related work}\label{related-work}

\subsubsection{2.1 Supply-chain delay prediction and pharmaceutical
logistics}\label{supply-chain-delay-prediction-and-pharmaceutical-logistics}

Machine learning has been applied broadly to supply-chain risk and
disruption prediction (Baryannis et al. 2019). Earlier pharmaceutical
work includes machine-learning lead-time forecasting (Biazon de Oliveira
et al. 2021), while more recent studies have examined pharmaceutical
disruption risk across the COVID-19 regime shift (Hupman, Zhang, and Li
2024) and global-health shipment delays using shipment and country-level
logistics indicators (Gali, Molavi, and Alavi 2025). Order-delay
classification is also represented in general supply-chain studies
(Thomas and Panicker 2023; Bassiouni et al. 2024). A recent analytical
review of delivery-delay prediction identifies continued emphasis on
predictive classification and highlights the importance of linking
predictions to decision support (Müller et al. 2025).

\subsubsection{2.2 Temporal evaluation, leakage, and label
maturity}\label{temporal-evaluation-leakage-and-label-maturity}

Temporally structured evaluation is necessary when observations are
dependent or non-stationary. Random cross-validation can understate
prospective error when data have temporal structure (Roberts et al.
2017), while leakage has been identified as a broader reproducibility
problem in machine-learning science (Kapoor and Narayanan 2023).
Time-series validation literature similarly emphasizes matching the
validation design to the deployment setting (Bergmeir and Benítez 2012).
Purging or embargo ideas have also been used in other time-dependent
domains to separate training observations from outcomes whose
information would not yet be available (López de Prado 2018). In the
present logistics setting, the embargo is tied directly to shipment
outcome resolution.

\subsubsection{2.3 Calibration and severity
uncertainty}\label{calibration-and-severity-uncertainty}

Probability calibration is distinct from discrimination. Platt scaling
applies a monotonic logistic mapping to model scores, while isotonic
regression provides a more flexible monotonic mapping (Platt 1999;
Zadrozny and Elkan 2002; Niculescu-Mizil and Caruana 2005). Reliability
is evaluated here with Brier score (Brier 1950) and a 10-bin expected
calibration error, with the latter treated as a descriptive calibration
metric rather than a sufficient stand-alone measure.

Conditional severity is modeled only among truly delayed shipments to
avoid a target distribution dominated by zero-delay observations.
Quantile regression (Koenker and Bassett 1978) provides lower, median,
and upper conditional estimates. CQR then adjusts learned quantile
intervals using held-out conformity scores (Romano, Patterson, and
Candès 2019). Standard conformal validity relies on exchangeability or
related conditions; covariate shift and broader non-exchangeability
require additional care (Tibshirani et al. 2019; Gibbs and Candès 2021;
Barber et al. 2023). We therefore distinguish formal
exchangeability-based guarantees from empirical future-period coverage
in this study.

\subsubsection{2.4 Prediction-to-decision integration and closest prior
art}\label{prediction-to-decision-integration-and-closest-prior-art}

Recent work narrows the novelty boundary substantially. In global-health
logistics, Pathak et al. (Pathak et al. 2025) proposed SCaLDR, a
two-stage framework for HIV-medicine shipment-delay occurrence and
duration. In broader logistics, Yang (Yang 2026) jointly modeled risk
classification and delay forecasting using a chronological design.
Faulkner et al. (Faulkner et al. 2026), a 2026 preprint, combines delay
classification, quantile duration modeling, and conformalized
uncertainty on a very large industrial shipment dataset. Conformal
prediction has also been propagated into capacity-constrained
container-terminal scheduling (Makhado, Sejeso, and Paepae 2026) and
transportation-delay uncertainty (Sadeek, Hanaoka, and Sugishita 2026).
Separately, delay classification has been coupled to supplier-selection
and order-allocation optimization (Zaghdoudi et al. 2024), and
risk-based container inspection has been formulated under explicit
operational resource constraints (Liang et al. 2026).

These studies rule out broad claims that the present work is the first
application of machine learning to pharmaceutical delays, the first use
of conformal prediction in logistics, or the first combination of delay
occurrence and duration. \textbf{To the best of our knowledge, this is
the first study to jointly evaluate, in pharmaceutical shipment
logistics, a leakage-aware temporally ordered delay-risk pipeline with
post-delivery embargoes, post-hoc probability calibration, conditional
delay-severity modeling, conformalized quantile uncertainty, and
capacity-constrained shipment prioritization.} This is a qualified
joint-combination priority claim based on literature reviewed through 31
August 2026, not a claim of algorithmic novelty.

\subsection{3. Data and prediction-time
formulation}\label{data-and-prediction-time-formulation}

\subsubsection{3.1 Dataset and outcome}\label{dataset-and-outcome}

The study uses the public USAID SCMS Delivery History Dataset and
associated program documentation (USAID SCMS Project 2015; USAID 2016).
The canonical file contains 10,324 completed shipment records across 42
recipient countries. There are 1,186 delayed records, corresponding to
an overall delay prevalence of 11.488\%. The commodity mix includes
antiretroviral pharmaceuticals, rapid diagnostic tests, antimalarial
products, and ancillary health-supply items. The dataset is historical
and observational; it does not contain randomized intervention outcomes.

For shipment i, the binary target is Y\_i = 1 when the delivered date
occurs after the scheduled delivery date and Y\_i = 0 otherwise.
Positive delay severity S\_i is the number of calendar days beyond the
scheduled date and is modeled only in the subset Y\_i = 1.

\subsubsection{3.2 Prediction-time feature
contract}\label{prediction-time-feature-contract}

The prediction anchor T\_pred represents the point at which the shipment
schedule is available and a prospective risk score could be issued. The
frozen research contract contains 39 pre-outcome features. Post-outcome
and target-derived fields - including Delay\_Days, Delay\_Flag,
Delivered to Client Date, and Delivery Recorded Date - are excluded.
Categorical missing values are represented explicitly; numeric
missingness is imputed using training-only statistics. The objective is
not merely to remove obvious target columns but to ensure that all
engineered features are computable at or before T\_pred.

\subsection{4. Temporal evaluation and leakage-control
protocol}\label{temporal-evaluation-and-leakage-control-protocol}

The primary evidence comes from five expanding-origin temporal
development folds covering a development cohort of 7,306 shipments,
including 1,125 delayed observations. A 90-day embargo is inserted
between the end of eligible training outcomes and each validation
window. The embargo acts as a label-maturity buffer so that shipments
whose final outcomes would not yet be resolved cannot contribute target
information to training. Figure 2 provides a schematic representation.

A later 717-row calibration buffer (103 delayed shipments) is used for
frozen post-hoc calibration and conformal adjustment where specified.
The secondary Locked Registry Evaluation Set contains 1,013 shipments
from 24 August 2014 through 24 August 2015, of which 61 are delayed
(6.02\%). This cohort was historically evaluated by the serving registry
before the research track and is therefore treated as a
\textbf{secondary locked registry benchmark}, not as a newly untouched
confirmatory holdout.

\textbf{Table 1. Evidence hierarchy and temporal cohorts.}

\begin{longtable}[]{@{}
  >{\raggedright\arraybackslash}p{(\columnwidth - 6\tabcolsep) * \real{0.2143}}
  >{\raggedleft\arraybackslash}p{(\columnwidth - 6\tabcolsep) * \real{0.2857}}
  >{\raggedleft\arraybackslash}p{(\columnwidth - 6\tabcolsep) * \real{0.2857}}
  >{\raggedright\arraybackslash}p{(\columnwidth - 6\tabcolsep) * \real{0.2143}}@{}}
\toprule\noalign{}
\begin{minipage}[b]{\linewidth}\raggedright
Cohort
\end{minipage} & \begin{minipage}[b]{\linewidth}\raggedleft
N
\end{minipage} & \begin{minipage}[b]{\linewidth}\raggedleft
Delayed
\end{minipage} & \begin{minipage}[b]{\linewidth}\raggedright
Scientific role
\end{minipage} \\
\midrule\noalign{}
\endhead
\bottomrule\noalign{}
\endlastfoot
Canonical SCMS population & 10,324 & 1,186 & Historical source
population \\
Temporal development cohort & 7,306 & 1,125 & Primary evidence: 5
expanding folds \\
Calibration buffer & 717 & 103 & Frozen probability/CQR calibration
buffer \\
Locked registry benchmark & 1,013 & 61 & Secondary replication/benchmark
evidence \\
\end{longtable}

\subsection{5. Delay-risk
classification}\label{delay-risk-classification}

Five tabular classifiers were compared: regularized logistic regression,
Random Forest, CatBoost, XGBoost, and LightGBM (Breiman 2001;
Prokhorenkova et al. 2018; Chen and Guestrin 2016; Ke et al. 2017).
PR-AUC is primary because delay is the minority class; ROC-AUC and Brier
score are secondary. Random 5-fold cross-validation is included only as
a diagnostic comparator. Model selection and paper claims rely on
temporal evaluation.

\textbf{Table 2. Random-split diagnostic versus expanding temporal
evaluation.}

\begin{longtable}[]{@{}
  >{\raggedright\arraybackslash}p{(\columnwidth - 10\tabcolsep) * \real{0.1304}}
  >{\raggedleft\arraybackslash}p{(\columnwidth - 10\tabcolsep) * \real{0.1739}}
  >{\raggedleft\arraybackslash}p{(\columnwidth - 10\tabcolsep) * \real{0.1739}}
  >{\raggedleft\arraybackslash}p{(\columnwidth - 10\tabcolsep) * \real{0.1739}}
  >{\raggedleft\arraybackslash}p{(\columnwidth - 10\tabcolsep) * \real{0.1739}}
  >{\raggedleft\arraybackslash}p{(\columnwidth - 10\tabcolsep) * \real{0.1739}}@{}}
\toprule\noalign{}
\begin{minipage}[b]{\linewidth}\raggedright
Model
\end{minipage} & \begin{minipage}[b]{\linewidth}\raggedleft
Random PR-AUC
\end{minipage} & \begin{minipage}[b]{\linewidth}\raggedleft
Temporal PR-AUC
\end{minipage} & \begin{minipage}[b]{\linewidth}\raggedleft
Relative inflation
\end{minipage} & \begin{minipage}[b]{\linewidth}\raggedleft
Temporal ROC-AUC
\end{minipage} & \begin{minipage}[b]{\linewidth}\raggedleft
Temporal Brier
\end{minipage} \\
\midrule\noalign{}
\endhead
\bottomrule\noalign{}
\endlastfoot
Logistic Regression & 0.3799 & 0.3011 & +26.2\% & 0.7132 & 0.2001 \\
Random Forest & 0.4780 & 0.3238 & +47.6\% & 0.7316 & 0.1605 \\
CatBoost & 0.5608 & 0.3164 & +77.3\% & 0.7401 & 0.1417 \\
XGBoost & 0.5591 & 0.2917 & +91.6\% & 0.7128 & 0.1395 \\
LightGBM & 0.5975 & 0.2992 & +99.7\% & 0.7277 & 0.1430 \\
\end{longtable}

Random splitting produced higher PR-AUC for every tested architecture.
Relative inflation ranged from 26.2\% for logistic regression to 99.7\%
for LightGBM, with Random Forest at 47.6\% and CatBoost at 77.3\%. We
interpret these differences as evidence that random splitting materially
overstates performance relative to the temporally ordered deployment
proxy. The experiment does not establish a single causal mechanism such
as entity memorization; the gap is consistent with temporal dependence
and changing distribution structure.

\subsection{6. Probability calibration}\label{probability-calibration}

Tree-ensemble scores were evaluated in raw form and after Platt or
isotonic calibration. Platt scaling was frozen for the operational
comparison because it improved probability reliability while preserving
continuous score ordering. Isotonic regression often improved ECE or
Brier score further but introduced ties that reduced PR-AUC in these
experiments.

\textbf{Table 3. Probability calibration on temporal development folds.}

\begin{longtable}[]{@{}llrrrr@{}}
\toprule\noalign{}
Model & Calibration & Brier & ECE (10 bins) & PR-AUC & ROC-AUC \\
\midrule\noalign{}
\endhead
\bottomrule\noalign{}
\endlastfoot
CatBoost & Raw & 0.1398 & 0.0850 & 0.2999 & 0.6916 \\
CatBoost & Platt & 0.1357 & 0.0807 & 0.2999 & 0.6916 \\
CatBoost & Isotonic & 0.1336 & 0.0760 & 0.2730 & 0.6854 \\
Random Forest & Raw & 0.1797 & 0.2051 & 0.3296 & 0.7269 \\
Random Forest & Platt & 0.1317 & 0.0866 & 0.3296 & 0.7269 \\
Random Forest & Isotonic & 0.1316 & 0.0774 & 0.2904 & 0.6985 \\
\end{longtable}

For CatBoost, Platt calibration reduced Brier score from 0.1398 to
0.1357 and ECE from 0.0850 to 0.0807 while leaving PR-AUC at 0.2999. For
Random Forest, the reduction was larger: Brier fell from 0.1797 to
0.1317 and ECE from 0.2051 to 0.0866 while PR-AUC remained 0.3296. No
hypothesis test of calibration improvement is claimed; the results are
descriptive quantitative comparisons.

\subsection{7. Conditional delay-severity
modeling}\label{conditional-delay-severity-modeling}

Severity is evaluated only on delayed shipments. A simple conditional
median provides a transparent baseline. Ridge regression tests a linear
feature-dependent model, while LightGBM quantile regressors estimate
asymmetric conditional quantiles used by the CQR layer. This separation
avoids allowing the large mass at zero delay to dominate a continuous
regression loss.

\textbf{Table 4. Conditional severity point error on temporal
development folds.}

\begin{longtable}[]{@{}lrrr@{}}
\toprule\noalign{}
Model & MAE, days & Median AE, days & q0.50 pinball \\
\midrule\noalign{}
\endhead
\bottomrule\noalign{}
\endlastfoot
Conditional median & 15.62 +/- 6.43 & 7.60 +/- 0.89 & 7.81 \\
LightGBM quantile q0.50 & 16.96 +/- 4.66 & 8.74 +/- 1.64 & 8.48 \\
Ridge regression & 23.76 +/- 3.92 & 15.01 +/- 0.54 & 11.88 \\
\end{longtable}

The conditional median achieved the lowest mean point error (15.62
days), compared with 16.96 days for the LightGBM median quantile and
23.76 days for Ridge. This negative result is important: the more
complex quantile model is not presented as a universally better point
predictor. Its principal role is to represent conditional asymmetry and
provide quantile bounds for conformalization.

\subsection{8. Conformalized quantile
regression}\label{conformalized-quantile-regression}

For nominal coverage 1-gamma, lower and upper LightGBM quantile models
estimate q\_(gamma/2)(x) and q\_(1-gamma/2)(x). On a disjoint
calibration set, the nonconformity score is E\_i =
max(q\_low(x\_i)-s\_i, s\_i-q\_high(x\_i)). The finite-sample adjustment
uses the empirical quantile of conformity scores with the frozen
higher-quantile rule. The final interval expands the learned lower and
upper quantiles by this adjustment.

Under exchangeability, split CQR provides finite-sample marginal
coverage guarantees (Romano, Patterson, and Candès 2019). Because the
evaluation here is explicitly temporal and may be non-exchangeable,
future-cohort coverage is interpreted as an \textbf{empirical
out-of-time audit}, not as a formal distribution-free guarantee under
arbitrary temporal shift (Tibshirani et al. 2019; Barber et al. 2023).

\textbf{Table 5. CQR coverage and sharpness.}

\begin{longtable}[]{@{}
  >{\raggedleft\arraybackslash}p{(\columnwidth - 12\tabcolsep) * \real{0.1429}}
  >{\raggedleft\arraybackslash}p{(\columnwidth - 12\tabcolsep) * \real{0.1429}}
  >{\raggedleft\arraybackslash}p{(\columnwidth - 12\tabcolsep) * \real{0.1429}}
  >{\raggedleft\arraybackslash}p{(\columnwidth - 12\tabcolsep) * \real{0.1429}}
  >{\raggedleft\arraybackslash}p{(\columnwidth - 12\tabcolsep) * \real{0.1429}}
  >{\raggedleft\arraybackslash}p{(\columnwidth - 12\tabcolsep) * \real{0.1429}}
  >{\raggedleft\arraybackslash}p{(\columnwidth - 12\tabcolsep) * \real{0.1429}}@{}}
\toprule\noalign{}
\begin{minipage}[b]{\linewidth}\raggedleft
Nominal
\end{minipage} & \begin{minipage}[b]{\linewidth}\raggedleft
Development coverage
\end{minipage} & \begin{minipage}[b]{\linewidth}\raggedleft
Dev. mean width (days)
\end{minipage} & \begin{minipage}[b]{\linewidth}\raggedleft
Dev. median width
\end{minipage} & \begin{minipage}[b]{\linewidth}\raggedleft
Locked coverage (61 delays)
\end{minipage} & \begin{minipage}[b]{\linewidth}\raggedleft
Exact 95\% CI
\end{minipage} & \begin{minipage}[b]{\linewidth}\raggedleft
Locked mean width
\end{minipage} \\
\midrule\noalign{}
\endhead
\bottomrule\noalign{}
\endlastfoot
80\% & 78.33\% +/- 8.37\% & 46.37 & 40.55 & 70.49\% (43/61) &
57.43-81.48\% & 41.04 \\
90\% & 85.76\% +/- 15.24\% & 110.33 & 106.73 & 91.80\% (56/61) &
81.90-97.28\% & 46.27 \\
95\% & 95.84\% +/- 6.02\% & 153.55 & 151.98 & 100.00\% (61/61) &
94.13-100.00\% & 61.74 \\
\end{longtable}

Development results show meaningful temporal variability, especially at
the 90\% level, where average interval width was 110.33 days with large
fold-to-fold variation. On the locked benchmark, the 80\% point estimate
undercovered (70.49\%), although the nominal 80\% level falls inside the
exact interval; the 90\% level achieved 91.80\% coverage; and the 95\%
level covered all 61 delayed shipments at the cost of wider intervals.
Reporting all three levels makes the coverage-sharpness trade-off
explicit rather than selecting a favorable level after evaluation.

\subsection{9. Secondary locked registry
benchmark}\label{secondary-locked-registry-benchmark}

The secondary benchmark is used to examine whether the frozen research
pipeline behaves consistently in the same historical period previously
seen by the serving registry. It is not a new globally unseen test set.
CatBoost is retained as the \textbf{deployment-aligned primary model},
while Random Forest is reported as the \textbf{development PR-AUC
sensitivity comparator}. Their roles were frozen before this benchmark
and are not redefined based on benchmark performance.

\textbf{Table 6. Locked registry benchmark classification results.}

\begin{longtable}[]{@{}
  >{\raggedright\arraybackslash}p{(\columnwidth - 20\tabcolsep) * \real{0.0714}}
  >{\raggedright\arraybackslash}p{(\columnwidth - 20\tabcolsep) * \real{0.0714}}
  >{\raggedleft\arraybackslash}p{(\columnwidth - 20\tabcolsep) * \real{0.0952}}
  >{\raggedleft\arraybackslash}p{(\columnwidth - 20\tabcolsep) * \real{0.0952}}
  >{\raggedleft\arraybackslash}p{(\columnwidth - 20\tabcolsep) * \real{0.0952}}
  >{\raggedleft\arraybackslash}p{(\columnwidth - 20\tabcolsep) * \real{0.0952}}
  >{\raggedleft\arraybackslash}p{(\columnwidth - 20\tabcolsep) * \real{0.0952}}
  >{\raggedleft\arraybackslash}p{(\columnwidth - 20\tabcolsep) * \real{0.0952}}
  >{\raggedleft\arraybackslash}p{(\columnwidth - 20\tabcolsep) * \real{0.0952}}
  >{\raggedleft\arraybackslash}p{(\columnwidth - 20\tabcolsep) * \real{0.0952}}
  >{\raggedleft\arraybackslash}p{(\columnwidth - 20\tabcolsep) * \real{0.0952}}@{}}
\toprule\noalign{}
\begin{minipage}[b]{\linewidth}\raggedright
Model
\end{minipage} & \begin{minipage}[b]{\linewidth}\raggedright
Role
\end{minipage} & \begin{minipage}[b]{\linewidth}\raggedleft
Threshold
\end{minipage} & \begin{minipage}[b]{\linewidth}\raggedleft
PR-AUC
\end{minipage} & \begin{minipage}[b]{\linewidth}\raggedleft
ROC-AUC
\end{minipage} & \begin{minipage}[b]{\linewidth}\raggedleft
Brier
\end{minipage} & \begin{minipage}[b]{\linewidth}\raggedleft
ECE
\end{minipage} & \begin{minipage}[b]{\linewidth}\raggedleft
Precision
\end{minipage} & \begin{minipage}[b]{\linewidth}\raggedleft
Recall
\end{minipage} & \begin{minipage}[b]{\linewidth}\raggedleft
F1
\end{minipage} & \begin{minipage}[b]{\linewidth}\raggedleft
Balanced accuracy
\end{minipage} \\
\midrule\noalign{}
\endhead
\bottomrule\noalign{}
\endlastfoot
CatBoost & Deployment-aligned primary & 0.1000 & 0.2709 & 0.7477 &
0.0527 & 0.0568 & 0.2174 & 0.7377 (45/61) & 0.3358 & 0.7899 \\
Random Forest & PR-AUC sensitivity comparator & 0.1050 & 0.3195 & 0.7898
& 0.0493 & 0.0314 & 0.2117 & 0.7705 (47/61) & 0.3322 & 0.8037 \\
\end{longtable}

Random Forest achieved higher benchmark PR-AUC (0.3195 versus 0.2709),
ROC-AUC, Brier score, ECE, recall, and balanced accuracy. CatBoost had
slightly higher F1 and precision at its frozen threshold. These results
do not support calling CatBoost the best predictive model; they support
transparent model-role separation and sensitivity analysis.

On the 61 delayed benchmark shipments, the conditional median achieved
MAE 9.05 days and median absolute error 7.00 days, whereas the LightGBM
quantile median had MAE 14.21 days but median absolute error 4.91 days.
Again, the evidence favors different models for different loss functions
rather than a universal severity winner.

\subsection{10. Capacity-constrained operational prioritization
{[}SIMULATED
SCENARIO{]}}\label{capacity-constrained-operational-prioritization-simulated-scenario}

The operational analysis assumes that a logistics team can inspect only
the top K\% of shipments. Three rankings are compared: calibrated risk
alone, calibrated risk multiplied by predicted median severity, and
calibrated risk multiplied by the upper-severity estimate. Outcomes are
retrospectively evaluated using observed delays; no real intervention
was applied, and no causal or financial-effect claim is made.

At K=1\%, risk-only ranking captured 1 of 15 high-severity delays, while
both severity-aware rankings captured 8 of 15. At K=5\%, the
corresponding counts were 2 of 15 versus 10 of 15. At K=10\%, both
severity-aware strategies captured 13 of 15 high-severity delays, but
risk x q50 captured more total delays than risk x q95. At K=20\%,
risk-only ranking captured all 15 high-severity delays and more total
delayed shipments than the other strategies. Thus, the evidence does not
support a universally superior uncertainty-aware ranking. Instead,
\textbf{risk, severity, and uncertainty generate different
prioritization trade-offs under constrained inspection capacity}.

The complete Top-K table is provided in the supplementary material and
every operational result is labeled \textbf{{[}SIMULATED SCENARIO{]}}.

\subsection{11. Discussion}\label{discussion}

\subsubsection{11.1 Random-split optimism is a logistics evaluation
problem}\label{random-split-optimism-is-a-logistics-evaluation-problem}

The largest and most consistent result is not that one classifier
dominates another, but that the evaluation design changes the apparent
level of performance. Every architecture showed higher PR-AUC under
random cross-validation than under expanding temporal evaluation. The
effect was model-dependent: 26.2\% for logistic regression and nearly
100\% for LightGBM. This reinforces the need to make prediction-time
information sets and outcome maturity explicit when evaluating logistics
forecasting systems.

The 90-day embargo is not proposed as a universally optimal constant. It
is a domain-specific label-maturity policy chosen to reduce the chance
that unresolved shipment outcomes leak across temporal boundaries. Other
logistics systems should set this gap from their own information-delay
characteristics.

\subsubsection{11.2 Calibration matters when probability enters a
decision
rule}\label{calibration-matters-when-probability-enters-a-decision-rule}

For resource allocation, the numeric scale of risk matters. Platt
scaling reduced Brier and ECE without changing continuous ordering,
making it an operationally convenient calibrator in this study. Isotonic
calibration sometimes attained lower calibration error but produced
coarser tied probabilities and lower PR-AUC. The implication is not that
Platt scaling is universally superior, but that calibrator choice should
reflect both reliability and the decision process consuming the scores.

\subsubsection{11.3 A simple severity baseline remains difficult to
beat}\label{a-simple-severity-baseline-remains-difficult-to-beat}

The conditional median's lower MAE is a useful caution against assuming
that model complexity automatically improves operational point
forecasts. Feature-dependent quantile models remain valuable because the
task is not only point prediction: asymmetric quantiles define
heterogeneous uncertainty bounds for CQR. The manuscript therefore
separates the point-estimation result from the uncertainty-estimation
role.

\subsubsection{11.4 Conformal uncertainty should be audited under
temporal
shift}\label{conformal-uncertainty-should-be-audited-under-temporal-shift}

The locked 90\% CQR result is encouraging in that period, but it should
not be generalized into a guarantee under shift. Development folds show
large variation in both coverage and width, and the 80\% locked point
estimate undercovers. The study therefore treats future coverage as an
empirical diagnostic. For live deployment, uncertainty quality would
require ongoing monitoring and governed recalibration rather than
permanent reliance on a frozen conformal layer.

\subsubsection{11.5 Prioritization depends on operational objective and
capacity}\label{prioritization-depends-on-operational-objective-and-capacity}

The simulated decision layer highlights an operational distinction often
hidden by aggregate classification metrics. Under tight review budgets,
incorporating severity strongly increased high-severity capture relative
to risk-only ranking. At larger budgets, risk-only ranking recovered
more total delayed shipments and eventually all high-severity delays.
This pattern is useful because it prevents the operational analysis from
collapsing into a single `best policy' headline. Organizations should
define whether the objective is broad delay capture, severe-delay
avoidance, delay-day capture, or another cost-sensitive criterion before
choosing a queue score.

\subsection{12. Limitations}\label{limitations}

First, the empirical evidence comes from one historical USAID SCMS
dataset. The multi-country nature of the data does not establish
generalizability to contemporary commercial, e-commerce, cold-chain, or
other supply chains. Second, all records are observational; no
randomized expediting or inspection interventions are available, so
operational triage is simulated and no causal effect or realized
financial savings are claimed. Third, the locked benchmark contains only
61 delayed shipments. A single covered or uncovered shipment therefore
changes empirical coverage by about 1.64 percentage points, motivating
exact binomial intervals.

Fourth, temporal prevalence and feature distributions vary, so direct
metric comparisons across periods must be contextualized. Brier score in
particular is prevalence-sensitive. Fifth, standard CQR validity
conditions do not automatically hold under temporal distribution shift.
The out-of-time coverage results are therefore empirical audits. Sixth,
the secondary locked registry benchmark was historically evaluated by
the serving registry before this research track and should not be
described as a newly untouched confirmatory holdout. Seventh, the
Level-2 priority statement is qualified by the literature reviewed
through 31 August 2026 and could be narrowed by newly indexed or
inaccessible prior work.

\subsection{13. Reproducibility and evidence
governance}\label{reproducibility-and-evidence-governance}

The research process freezes the dataset, split logic, model roles,
thresholds, calibration policy, and benchmark protocol before the
secondary benchmark is read. The canonical dataset SHA-256 is
\texttt{918b992dd3e8d4b64d2a727b2c4ea607603d0c58f19484e73f7b78528c6a8673}.
The final evaluation freeze SHA-256 is
\texttt{631F1FA4241FBE697B46D9658B8720D2CE13CFB5A9F65E32C7532F8A1F6CD21A},
and the locked benchmark manifest SHA-256 is
\texttt{62BE80CE9BC977767BC983EDABF61DAA23609A1801251F4E93CFA9693EBDD9AB}.
The frozen integrity suite reported 26 of 26 tests passing.

For double-anonymized review, repository identity and author-linked URLs
are omitted from this manuscript. A blinded reproducibility package can
be supplied during review, with public repository details restored after
the review process in accordance with journal policy.

\subsection{14. Conclusion}\label{conclusion}

This study evaluates a logistics delay-intelligence pipeline under
temporally ordered conditions rather than treating a random split as a
proxy for deployment. Across five classifier families, random splitting
materially overstated PR-AUC. Probability calibration improved
reliability without requiring a new classifier, while a simple
conditional median remained competitive for severity point estimation.
CQR added an explicit uncertainty layer whose future-period performance
was evaluated empirically rather than assumed to remain valid under
shift. Finally, simulated Top-K prioritization showed that risk,
severity, and uncertainty support different operational objectives
depending on review capacity.

The main contribution is therefore methodological integration and
evaluation discipline rather than a new learning algorithm. For
pharmaceutical shipment logistics, the results support a workflow in
which prediction-time leakage control, temporally ordered validation,
calibrated risk, conditional severity, uncertainty auditing, and
explicit operational capacity are treated as one connected
decision-support problem.

\subsection{Data availability}\label{data-availability}

The study uses the publicly released USAID SCMS Delivery History
Dataset. The exact source citation and data-access route are provided in
the references. A checksum of the canonical analysis file is reported in
the reproducibility section.

\subsection{Code availability}\label{code-availability}

Code and frozen research artifacts are maintained in a
version-controlled research repository. Repository identity is
\textbf{{[}BLINDED FOR DOUBLE-ANONYMIZED REVIEW{]}} in this submission
version. The final accepted manuscript will provide the public
repository and reproducibility instructions, subject to journal policy.

\subsection{Ethics and responsible
use}\label{ethics-and-responsible-use}

The analysis uses retrospective shipment transactions and is intended
for logistics decision support. Predictions should not be used as
automatic punitive assessments of suppliers or personnel. Operational
rankings require human review, and simulated decision outcomes should
not be interpreted as observed intervention effects.

\subsection{Declaration of generative AI and AI-assisted technologies in
the manuscript preparation
process}\label{declaration-of-generative-ai-and-ai-assisted-technologies-in-the-manuscript-preparation-process}

During the preparation of this work, the author used OpenAI ChatGPT and
Google Gemini to support literature-search organization, manuscript
structuring, language refinement, bibliographic integrity checks, and
submission-readiness review. After using these tools/services, the author
reviewed and edited the content as needed, independently verified
quantitative results and references against frozen artifacts and primary
sources, and takes full responsibility for the content of the publication.

\subsection{References}\label{references}

\begin{enumerate}
\def\labelenumi{\arabic{enumi}.}
\tightlist
\item
  Müller, N., Burggräf, P., Steinberg, F., Sauer, C.R., Schütz, M.
  (2025). An analytical review of predictive methods for delivery delays
  in supply chains. Supply Chain Analytics 11, 100130.
  https://doi.org/10.1016/j.sca.2025.100130
\item
  Biazon de Oliveira, M., Zucchi, G., Lippi, M., Cordeiro, D.F., Rosa da
  Silva, N., Iori, M. (2021). Lead Time Forecasting with Machine
  Learning Techniques for a Pharmaceutical Supply Chain. ICEIS, 634-641.
  https://doi.org/10.5220/0010434406340641
\item
  Pathak, N., Keshari, S., Rai, S., Saksham, K., Raj, P. (2025).
  Predicting Global Supply Chain Delays for HIV Medicines Using SCaLDR:
  A Two-stage ML Framework. IEEE CICT.
  https://doi.org/10.1109/CICT67193.2025.11399211
\item
  Yang, X. (2026). O²RDL-net for joint risk classification and delay
  forecasting in logistics systems using interaction amplified deep.
  Scientific Reports 16, 24683.
  https://doi.org/10.1038/s41598-026-55703-6
\item
  Faulkner, S., Zandehshahvar, R., Eghbal Akhlaghi, V., Ouellet, S.,
  Jordan, C., Van Hentenryck, P. (2026). Uncertainty-Aware Delivery
  Delay Duration Prediction via Multi-Task Deep Learning.
  arXiv:2602.20271. Preprint.
\item
  Makhado, N., Sejeso, M., Paepae, T. (2026). Conformal prediction-based
  sequential scheduling for container terminal operations under
  uncertainty. Operations Research Perspectives 17, 100409.
  https://doi.org/10.1016/j.orp.2026.100409
\item
  Liang, X., Fan, S., Li, H., Goerlandt, F., Yang, Z. (2026). Freeports
  under the lens: securing container supply chains with a risk-based
  inspection framework. Transportation Research Part E 208, 104658.
  https://doi.org/10.1016/j.tre.2025.104658
\item
  Zaghdoudi, M.A., Hajri-Gabouj, S., Ghezail, F., Darmoul, S., Varnier,
  C., Zerhouni, N. (2024). Collaborative and integrated data-driven
  delay prediction and supplier selection optimization: A case study in
  a furniture industry. Computers \& Industrial Engineering 197, 110590.
  https://doi.org/10.1016/j.cie.2024.110590
\item
  Hupman, A.C., Zhang, J., Li, H. (2024). Predicting pharmaceutical
  supply chain disruptions before and during the COVID-19 pandemic. Risk
  Analysis 44(12), 2797-2811. https://doi.org/10.1111/risa.17453
\item
  Gali, J.S., Molavi, N., Alavi, S. (2025). Predicting Global Healthcare
  Supply Chain Delays: A Machine Learning Approach Leveraging
  Country-level Logistics Metrics. Journal of International Technology
  and Information Management 34(1), Article 3.
  https://doi.org/10.58729/1941-6679.1634
\item
  Sadeek, S.N., Hanaoka, S., Sugishita, K. (2026). Uncertainty
  quantification of departure delay considering network properties and
  conformal prediction framework. Journal of Air Transport Management
  136, 103037. https://doi.org/10.1016/j.jairtraman.2026.103037
\item
  Bassiouni, M.M., Chakrabortty, R.K., Sallam, K.M., Hussain, O.K.
  (2024). Deep learning approaches to identify order status in a complex
  supply chain. Expert Systems with Applications 250, 123947.
  https://doi.org/10.1016/j.eswa.2024.123947
\item
  Thomas, A., Panicker, V.V. (2023). Application of Machine Learning
  Algorithms for Order Delivery Delay Prediction in Supply Chain
  Disruption Management. In Intelligent Manufacturing Systems in
  Industry 4.0, 491-500. https://doi.org/10.1007/978-981-99-1665-8\_42
\item
  Kapoor, S., Narayanan, A. (2023). Leakage and the reproducibility
  crisis in machine-learning-based science. Patterns 4(9), 100804.
  https://doi.org/10.1016/j.patter.2023.100804
\item
  Roberts, D.R. et al.~(2017). Cross-validation strategies for data with
  temporal, spatial or phylogenetic structure. Ecography 40(8), 913-929.
  https://doi.org/10.1111/ecog.02881
\item
  Niculescu-Mizil, A., Caruana, R. (2005). Predicting good probabilities
  with supervised learning. ICML, 625-632.
  https://doi.org/10.1145/1102351.1102430
\item
  Romano, Y., Patterson, E., Candès, E. (2019). Conformalized Quantile
  Regression. NeurIPS 32, 3543-3553.
\item
  Barber, R.F., Candès, E.J., Ramdas, A., Tibshirani, R.J. (2023).
  Conformal prediction beyond exchangeability. Annals of Statistics
  51(2), 816-845. https://doi.org/10.1214/23-AOS2276
\item
  Yadav, P. (2015). Health Product Supply Chains in Developing
  Countries. Health Systems \& Reform 1(2), 142-154.
  https://doi.org/10.4161/23288604.2014.968005
\item
  Vledder, M., Friedman, J., Sjöblom, M., Brown, T., Yadav, P. (2019).
  Improving Supply Chain for Essential Drugs in Low-Income Countries.
  Health Systems \& Reform 5(2), 158-177.
  https://doi.org/10.1080/23288604.2019.1596050
\item
  Baryannis, G., Validi, S., Dani, S., Antoniou, G. (2019). Supply chain
  risk management and artificial intelligence: state of the art and
  future research directions. International Journal of Production
  Research 57(7), 2179-2202.
  https://doi.org/10.1080/00207543.2018.1530476
\item
  Guo, C., Pleiss, G., Sun, Y., Weinberger, K.Q. (2017). On Calibration
  of Modern Neural Networks. ICML 70, 1321-1330.
\item
  Koenker, R., Bassett, G. (1978). Regression Quantiles. Econometrica
  46(1), 33-50. https://doi.org/10.2307/1913643
\item
  Bertsimas, D., Kallus, N. (2020). From Predictive to Prescriptive
  Analytics. Management Science 66(3), 1025-1044.
  https://doi.org/10.1287/mnsc.2018.3253
\item
  USAID SCMS Project (2015). Supply Chain Management System: Final
  Program Report (2005-2015). U.S. Agency for International Development
  and Partnership for Supply Chain Management.
\end{enumerate}

\phantomsection\label{refs}
\begin{CSLReferences}{1}{0}
\bibitem[\citeproctext]{ref-angelopoulos2023gentle}
Angelopoulos, Anastasios N., and Stephen Bates. 2023. {``Conformal
Prediction: A Gentle Introduction.''} \emph{Foundations and Trends in
Machine Learning} 16 (4): 494--591.
\url{https://doi.org/10.1561/2200000101}.

\bibitem[\citeproctext]{ref-barber2023beyond}
Barber, Rina Foygel, Emmanuel J. Candès, Aaditya Ramdas, and Ryan J.
Tibshirani. 2023. {``Conformal Prediction Beyond Exchangeability.''}
\emph{The Annals of Statistics} 51 (2): 816--45.
\url{https://doi.org/10.1214/23-AOS2276}.

\bibitem[\citeproctext]{ref-baryannis2019supply}
Baryannis, George, Sahar Validi, Samir Dani, and Grigoris Antoniou.
2019. {``Supply Chain Risk Management and Artificial Intelligence: State
of the Art and Future Research Directions.''} \emph{International
Journal of Production Research} 57 (7): 2179--2202.
\url{https://doi.org/10.1080/00207543.2018.1530476}.

\bibitem[\citeproctext]{ref-bassiouni2024deep}
Bassiouni, Mahmoud M., Ripon K. Chakrabortty, Karam M. Sallam, and Omar
K. Hussain. 2024. {``Deep Learning Approaches to Identify Order Status
in a Complex Supply Chain.''} \emph{Expert Systems with Applications}
250: 123947. \url{https://doi.org/10.1016/j.eswa.2024.123947}.

\bibitem[\citeproctext]{ref-bergmeir2012use}
Bergmeir, Christoph, and José M. Benítez. 2012. {``On the Use of
Cross-Validation for Time Series Predictor Evaluation.''}
\emph{Information Sciences} 191: 192--213.
\url{https://doi.org/10.1016/j.ins.2011.12.028}.

\bibitem[\citeproctext]{ref-bertsimas2020predictive}
Bertsimas, Dimitris, and Nathan Kallus. 2020. {``From Predictive to
Prescriptive Analytics.''} \emph{Management Science} 66 (3): 1025--44.
\url{https://doi.org/10.1287/mnsc.2018.3253}.

\bibitem[\citeproctext]{ref-oliveira2021lead}
Biazon de Oliveira, Maiza, Giorgio Zucchi, Marco Lippi, Douglas Farias
Cordeiro, Núbia Rosa da Silva, and Manuel Iori. 2021. {``Lead Time
Forecasting with Machine Learning Techniques for a Pharmaceutical Supply
Chain.''} In \emph{Proceedings of the 23rd International Conference on
Enterprise Information Systems (ICEIS)}, 1:634--41. SciTePress.
\url{https://doi.org/10.5220/0010434406340641}.

\bibitem[\citeproctext]{ref-breiman2001random}
Breiman, Leo. 2001. {``Random Forests.''} \emph{Machine Learning} 45
(1): 5--32. \url{https://doi.org/10.1023/A:1010933404324}.

\bibitem[\citeproctext]{ref-brier1950verification}
Brier, Glenn W. 1950. {``Verification of Forecasts Expressed in Terms of
Probability.''} \emph{Monthly Weather Review} 78 (1): 1--3.
\url{https://doi.org/10.1175/1520-0493(1950)078\%3C0001:VOFEIT\%3E2.0.CO;2}.

\bibitem[\citeproctext]{ref-chen2016xgboost}
Chen, Tianqi, and Carlos Guestrin. 2016. {``XGBoost: A Scalable Tree
Boosting System.''} In \emph{Proceedings of the 22nd ACM SIGKDD
International Conference on Knowledge Discovery and Data Mining},
785--94. \url{https://doi.org/10.1145/2939672.2939785}.

\bibitem[\citeproctext]{ref-faulkner2026uncertainty}
Faulkner, Stefan, Reza Zandehshahvar, Vahid Eghbal Akhlaghi, Sebastien
Ouellet, Carsten Jordan, and Pascal Van Hentenryck. 2026.
{``Uncertainty-Aware Delivery Delay Duration Prediction via Multi-Task
Deep Learning.''} \url{https://arxiv.org/abs/2602.20271}.

\bibitem[\citeproctext]{ref-gali2025predicting}
Gali, Jeevan Sai, Nima Molavi, and Sepideh Alavi. 2025. {``Predicting
Global Healthcare Supply Chain Delays: A Machine Learning Approach
Leveraging Country-Level Logistics Metrics.''} \emph{Journal of
International Technology and Information Management} 34 (1).
\url{https://doi.org/10.58729/1941-6679.1634}.

\bibitem[\citeproctext]{ref-gibbs2021adaptive}
Gibbs, Isaac, and Emmanuel Candès. 2021. {``Adaptive Conformal Inference
Under Distribution Shift.''} In \emph{Advances in Neural Information
Processing Systems}, 34:1660--72.

\bibitem[\citeproctext]{ref-guo2017calibration}
Guo, Chuan, Geoff Pleiss, Yu Sun, and Kilian Q. Weinberger. 2017. {``On
Calibration of Modern Neural Networks.''} In \emph{Proceedings of the
34th International Conference on Machine Learning}, 70:1321--30. PMLR.

\bibitem[\citeproctext]{ref-hupman2024predicting}
Hupman, Andrea C., Juan Zhang, and Haitao Li. 2024. {``Predicting
Pharmaceutical Supply Chain Disruptions Before and During the COVID-19
Pandemic.''} \emph{Risk Analysis} 44 (12): 2797--2811.
\url{https://doi.org/10.1111/risa.17453}.

\bibitem[\citeproctext]{ref-kapoor2023leakage}
Kapoor, Sayash, and Arvind Narayanan. 2023. {``Leakage and the
Reproducibility Crisis in Machine-Learning-Based Science.''}
\emph{Patterns} 4 (9): 100804.
\url{https://doi.org/10.1016/j.patter.2023.100804}.

\bibitem[\citeproctext]{ref-ke2017lightgbm}
Ke, Guolin, Qi Meng, Thomas Finley, Taifeng Wang, Wei Chen, Weidong Ma,
Qiwei Ye, and Tie-Yan Liu. 2017. {``LightGBM: A Highly Efficient
Gradient Boosting Decision Tree.''} In \emph{Advances in Neural
Information Processing Systems}, 30:3146--54.

\bibitem[\citeproctext]{ref-koenker1978regression}
Koenker, Roger, and Gilbert Bassett. 1978. {``Regression Quantiles.''}
\emph{Econometrica} 46 (1): 33--50.
\url{https://doi.org/10.2307/1913643}.

\bibitem[\citeproctext]{ref-liang2026freeports}
Liang, Xinrui, Shiqi Fan, Huanhuan Li, Floris Goerlandt, and Zaili Yang.
2026. {``Freeports Under the Lens: Securing Container Supply Chains with
a Risk-Based Inspection Framework.''} \emph{Transportation Research Part
E: Logistics and Transportation Review} 208: 104658.
\url{https://doi.org/10.1016/j.tre.2025.104658}.

\bibitem[\citeproctext]{ref-deprado2018advances}
López de Prado, Marcos. 2018. \emph{Advances in Financial Machine
Learning}. John Wiley \& Sons.

\bibitem[\citeproctext]{ref-makhado2026conformal}
Makhado, Ndifelani, Matthews Sejeso, and Thulane Paepae. 2026.
{``Conformal Prediction-Based Sequential Scheduling for Container
Terminal Operations Under Uncertainty.''} \emph{Operations Research
Perspectives} 17: 100409.
\url{https://doi.org/10.1016/j.orp.2026.100409}.

\bibitem[\citeproctext]{ref-muller2025analytical}
Müller, Norman, Peter Burggräf, Fabian Steinberg, Carl René Sauer, and
Maximilian Schütz. 2025. {``An Analytical Review of Predictive Methods
for Delivery Delays in Supply Chains.''} \emph{Supply Chain Analytics}
11: 100130. \url{https://doi.org/10.1016/j.sca.2025.100130}.

\bibitem[\citeproctext]{ref-niculescu2005predicting}
Niculescu-Mizil, Alexandru, and Rich Caruana. 2005. {``Predicting Good
Probabilities with Supervised Learning.''} In \emph{Proceedings of the
22nd International Conference on Machine Learning}, 625--32.
\url{https://doi.org/10.1145/1102351.1102430}.

\bibitem[\citeproctext]{ref-pathak2025predicting}
Pathak, N., S. Keshari, S. Rai, K. Saksham, and P. Raj. 2025.
{``Predicting Global Supply Chain Delays for HIV Medicines Using SCaLDR:
A Two-Stage ML Framework.''} In \emph{2025 IEEE 9th International
Conference on Information and Communication Technology (CICT)}. IEEE.
\url{https://doi.org/10.1109/CICT67193.2025.11399211}.

\bibitem[\citeproctext]{ref-platt1999probabilistic}
Platt, John. 1999. {``Probabilistic Outputs for Support Vector Machines
and Comparisons to Regularized Likelihood Methods.''} In \emph{Advances
in Large Margin Classifiers}, 61--74. MIT Press.

\bibitem[\citeproctext]{ref-prokhorenkova2018catboost}
Prokhorenkova, Liudmila, Gleb Gusev, Aleksandr Vorobev, Anna Veronika
Dorogush, and Andrey Gulin. 2018. {``CatBoost: Unbiased Boosting with
Categorical Features.''} In \emph{Advances in Neural Information
Processing Systems}, 31:6638--48.

\bibitem[\citeproctext]{ref-roberts2017crossvalidation}
Roberts, David R., Volker Bahn, Simone Ciuti, Mark S. Boyce, Jane Elith,
Gurutzeta Guillera-Arroita, Severin Hauenstein, et al. 2017.
{``Cross-Validation Strategies for Data with Temporal, Spatial or
Phylogenetic Structure.''} \emph{Ecography} 40 (8): 913--29.
\url{https://doi.org/10.1111/ecog.02881}.

\bibitem[\citeproctext]{ref-romano2019cqr}
Romano, Yaniv, Evan Patterson, and Emmanuel Candès. 2019.
{``Conformalized Quantile Regression.''} In \emph{Advances in Neural
Information Processing Systems}, 32:3543--53.

\bibitem[\citeproctext]{ref-sadeek2026uncertainty}
Sadeek, Soumik Nafis, Shinya Hanaoka, and Kashin Sugishita. 2026.
{``Uncertainty Quantification of Departure Delay Considering Network
Properties and Conformal Prediction Framework.''} \emph{Journal of Air
Transport Management} 136: 103037.
\url{https://doi.org/10.1016/j.jairtraman.2026.103037}.

\bibitem[\citeproctext]{ref-thomas2023application}
Thomas, Arun, and Vinay V. Panicker. 2023. {``Application of Machine
Learning Algorithms for Order Delivery Delay Prediction in Supply Chain
Disruption Management.''} In \emph{Intelligent Manufacturing Systems in
Industry 4.0: Select Proceedings of IPDIMS 2022}, 491--500. Springer
Nature Singapore. \url{https://doi.org/10.1007/978-981-99-1665-8_42}.

\bibitem[\citeproctext]{ref-tibshirani2019covariate}
Tibshirani, Ryan J., Rina Foygel Barber, Emmanuel Candès, and Aaditya
Ramdas. 2019. {``Conformal Prediction Under Covariate Shift.''} In
\emph{Advances in Neural Information Processing Systems}, 32:2530--40.

\bibitem[\citeproctext]{ref-scmsdataset2016}
USAID. 2016. {``USAID Supply Chain Management System (SCMS) Delivery
History Dataset.''} U.S. President's Emergency Plan for AIDS Relief
(PEPFAR) / USAID Data Repository.

\bibitem[\citeproctext]{ref-usaid2015scms}
USAID SCMS Project. 2015. {``Supply Chain Management System: Final
Program Report (2005--2015).''} Washington, DC: U.S. Agency for
International Development; Partnership for Supply Chain Management.

\bibitem[\citeproctext]{ref-vledder2019improving}
Vledder, Monique, Jed Friedman, Mirja Sjöblom, Timothy Brown, and
Prashant Yadav. 2019. {``Improving Supply Chain for Essential Drugs in
Low-Income Countries: Results from a Large Scale Randomized Experiment
in Zambia.''} \emph{Health Systems \& Reform} 5 (2): 158--77.
\url{https://doi.org/10.1080/23288604.2019.1596050}.

\bibitem[\citeproctext]{ref-yadav2015health}
Yadav, Prashant. 2015. {``Health Product Supply Chains in Developing
Countries: Diagnosis of the Root Causes of Underperformance and an
Agenda for Reform.''} \emph{Health Systems \& Reform} 1 (2): 142--54.
\url{https://doi.org/10.4161/23288604.2014.968005}.

\bibitem[\citeproctext]{ref-yang2026o2rdl}
Yang, Xiao. 2026. {``O\(^2\)RDL-Net for Joint Risk Classification and
Delay Forecasting in Logistics Systems Using Interaction Amplified
Deep.''} \emph{Scientific Reports} 16: 24683.
\url{https://doi.org/10.1038/s41598-026-55703-6}.

\bibitem[\citeproctext]{ref-zadrozny2002transforming}
Zadrozny, Bianca, and Charles Elkan. 2002. {``Transforming Classifier
Scores into Accurate Multiclass Probability Estimates.''} In
\emph{Proceedings of the Eighth ACM SIGKDD International Conference on
Knowledge Discovery and Data Mining}, 694--99.
\url{https://doi.org/10.1145/775047.775151}.

\bibitem[\citeproctext]{ref-zaghdoudi2024collaborative}
Zaghdoudi, Mohamed Aziz, Sonia Hajri-Gabouj, Feiza Ghezail, Saber
Darmoul, Christophe Varnier, and Noureddine Zerhouni. 2024.
{``Collaborative and Integrated Data-Driven Delay Prediction and
Supplier Selection Optimization: A Case Study in a Furniture
Industry.''} \emph{Computers \& Industrial Engineering} 197: 110590.
\url{https://doi.org/10.1016/j.cie.2024.110590}.

\end{CSLReferences}

\end{document}
